\documentclass[10pt,a4paper]{amsart}
\usepackage[margin=19mm]{geometry}
\usepackage{amsmath,amssymb,mathtools}
\usepackage[scr=rsfs,cal=euler]{mathalfa}
\usepackage{xfrac}
\usepackage[unicode,hidelinks,bookmarks=false]{hyperref}
\newcommand{\ie}{i.e.\ }
\newcommand{\cf}{cf.\ }
\newcommand{\resp}{respectively\ }
\newcommand{\Subsection}{\subsection}
\providecommand{\nobreakdash}{\nobreak\hbox{-}}
\setlength{\emergencystretch}{2em}
\multlinegap0pt
\numberwithin{equation}{section}
\newtheorem{theorem}{Theorem}[section]
\newtheorem{lemma}[theorem]{Lemma}
\newtheorem{definition}[theorem]{Definition}
\newtheorem{remark}[theorem]{Remark}
\allowdisplaybreaks
\title[Nonlinear stability of inviscid magnetic B\'enard fluids: original Theorem 1.1]{Nonlinear stability and optimal decay rates of inviscid magnetic B\'enard fluids: the original Theorem 1.1 and its complete original proof}
\author{Hao Liu, Ruixin Zeng, Chengrong Zhang}
\date{Source: Electronic Journal of Differential Equations 2026, No. 17, pp. 1--15; DOI 10.58997/ejde.2026.17; publisher version}
\begin{document}
\maketitle
\noindent\emph{Electronic Journal of Differential Equations}, Vol. 2026 (2026), No. 17, pp. 1--15.\newline
ISSN: 1072-6691. DOI \href{https://doi.org/10.58997/ejde.2026.17}{10.58997/ejde.2026.17}. This work is licensed under a CC BY 4.0 license.\par\smallskip
\noindent\textit{Editorial scope.} Counted claim: exactly one, the article's original Theorem 1.1 as printed, the paper's main nonlinear-stability theorem for the two-dimensional inviscid magnetic B\'enard system in homogeneous Besov spaces (source0.tex lines 164--190, the single primary statement), together with its complete original proof (source0.tex lines 887--908, the single primary proof), which is the author's own closing argument of the original Section 3. That argument is complete only together with the two frequency-localized a priori lemmas it invokes, and both are therefore retained verbatim and uncounted as the necessary complete core: original Lemma 3.1 (statement 521--526) with its complete original proof (528--694), and original Lemma 3.2 (statement 699--709) with its complete original proof (711--885). Also retained verbatim and uncounted as context: the whole of the original Section 1, Introduction (62--330), which carries the magnetic B\'enard system (1.1), the equilibrium state, the perturbation system (1.2), the initial conditions, the low-frequency and high-frequency notation, the energy functionals and the statements of the paper's Theorems 1.4 and 1.6 together with the author's outline of the proof; and the whole of the original Section 2, Preliminaries (333--498), namely the Littlewood--Paley decomposition, Definition 2.1 of the homogeneous Besov space, Lemma 2.2 (Bernstein inequalities), Lemma 2.3 (embedding, interpolation and derivative properties), Definition 2.4 of the mixed time-space space, Remark 2.5 and the product and commutator Lemmas 2.6 and 2.7; and the opening of the original Section 3 (500--519) with the energy functionals (3.1). No content line of the original Sections 1 to 3 is dropped, so every printed equation, theorem, lemma, definition and remark number of the delivered text coincides with the published version: the counted theorem is printed as Theorem 1.1, its a priori estimate as (1.3), the non-degeneracy condition as (1.4), the solution classes as (1.5), and the two support lemmas as Lemma 3.1 and Lemma 3.2. Excluded from this unit: the proof of Theorem 1.4 and the whole of the original Section 4 (910--1129), so the statements of Theorem 1.4 and Theorem 1.6 are carried as uncounted context whose proofs lie outside the unit; also excluded are the acknowledgements and the printed bibliography, whose entries are transcribed into the delivery with their published numbers. Printed slips are kept literal and no mathematics has been written, repaired or re-derived: the source prints \emph{Interested reads are referred to} (line 506) and \emph{we need to proof that} (line 889); it prints \emph{Berstein's inequality} for Bernstein's inequality (lines 372, 637, 652 and 765); the cross term displayed in (3.8) and again in (3.16) carries the coefficient nu on the term that pairs the Laplacian of the dyadic block of the magnetic field with the vertical derivative of the velocity, although the magnetic equation carries mu (lines 604 and 762); the first line of the displayed identity (3.6) prints the two norms without their exponents (line 581); the fifth auxiliary operator is printed with the factor (partial t minus nu) where the fourth carries (partial t plus nu) (line 306); the identity preceding (3.16) prints a malformed dyadic operator, a dot over a Delta followed by a braced Delta with subscript j (line 716); the integrated estimate (3.11) is printed for the un-squared energy functional while the display immediately before it is written for its square (lines 655--665); and the sentence \emph{This article paper is organized as follows} repeats its subject (line 319). Article-level slips outside the excerpt, left untouched: the article's own printed header prints Vol. 2025 (2025), No. 17 and DOI 10.58997/ejde.2025.17 with a 2025 copyright year (source0.tex lines 8--10 and page 1 of the publisher PDF) although the article is volume 2026, No. 17 under DOI 10.58997/ejde.2026.17, so the header line added by this editorial prelude carries the volume and DOI of record; and the published abstract breaks off in the fragment \emph{Furthermore, under appropriate additional conditions on low-frequency part.} (source0.tex lines 44--45).
 \section{Introduction}\label{sect1}

Magnetic B\'enard fluids describe the dynamic phenomenon of the interaction between velocity
field and magnetic field in conductive fluids, which plays a crucial role in thermal convection problems \cite{MR3363411,MR44996}.
In this article, we study the stability of 2D magnetic B\'enard fluids, described by the following system:
\begin{equation}\label{1.1}
\begin{gathered}
 u_t+ u\cdot \nabla u+\nabla P-\eta\Delta u= (0,\vartheta)^\mathrm{T}+ B\cdot \nabla B,\\
 B_t+u\cdot\nabla B-\mu\Delta B=B\cdot\nabla u,\\
 \vartheta_t+u\cdot \nabla \vartheta-\kappa\Delta\vartheta=0,\\
 \operatorname{div} u=\operatorname{div} B=0,
\end{gathered}
\end{equation}
where $u:=u(x,t)$, $B:=B(x,t)$, $P:=P(x,t)$ and $\vartheta:=\vartheta(x,t)$ represent the velocity field,
magnetic field, pressure and temperature of the fluid, respectively. The constants $\eta$, $\mu$ and
$\kappa$ denote the coefficients of viscosity, magnetic diffusivity and thermometric conductivity, respectively. The term $(0,\vartheta)^\mathrm{T}$ represents the buoyancy force acting due to the temperature variation.


The magnetic B\'enard fluid is a classical model for studying heat convection in the presence of
a magnetic field, and it plays a crucial role in Rayleigh--B\'enard convection. This system has broad
applications in both physics and geophysics \cite{MR53708,MR1965452}.
Significant progress has been made in understanding the global stability and large-time behavior of
solution for viscous, thermally and electrically conducting fluids, as demonstrated 
in \cite{MR3710679, MR64593,MR1107136,MR4134456,MR4505194,MR44301,MR3887133}.
In particular, the well-posedness of the system under partial dissipation has been extensively
studied in various settings, such as in 
\cite{MR4484087, MR4046053, MR3412279, MR4500250, MR3980230, MR4505194, 
MR4790886, MR4500891, MR3977115}.

When $\vartheta=0$, equation \eqref{1.1} reduces to the classical MHD system.
The global well-posedness of this system has been widely studied \cite{MR2418123,MR346289,MR716200}, particularly for the inviscid and non-resistive case
 $(\eta=\mu=0)$, as discussed in \cite{MR920153,MR3780143,MR3738384}.
 Initiated by Lin and Zhang \cite{MR3168121}, many papers
\cite{MR3666563, MR3780143, MR3851055, MR4339533, MR4641656,MR3377532, MR3210038, MR3678491,
MR3448784}
have focused on the global well-posedness of the MHD system with dissipation but without magnetic resistivity $(\eta>0,\mu=0)$.
For the inviscous and resistive MHD equations $(\eta=0,\mu>0)$, the global regularity problem has attracted attention. Under the certain symmetry assumptions, Zhou--Zhu \cite{MR3843627} studied
the global existence of classical solutions near an equilibrium state in the periodic domain.
Utilizing the assumption that $\int_{\mathbb{T}^2}b_0{\rm d}x=0$, Wei--Zhang \cite{MR4199961}
proved the global well-posedness of the system without a non-trivial background magnetic field in
$H^4$. Later, Ye--Yin \cite{MR4496608} conducted small solution with lower regularity in $H^s (s>2)$. Assuming the initial magnetic field is close to a background magnetic field satisfying the Diophantine condition, Zhai \cite{MR4624442} and Chen--Zhang--Zhou \cite{MR4372923} proved the global well-posedness in the periodic domain $\mathbb{T}^2$ and $\mathbb{T}^3$, respectively. Based on these researches, all globally well-posed results have been established in bounded domains. However, constructing a global solution for this inviscid system, even with small initial data, in the whole space remains an extremely challenging open problem.

 The goal of this article is to address this issue within the context of Besov spaces.
When the inviscid fluid is coupled with the temperature through the magnetic B\'enard system
\eqref{1.2}, which involves a thermal damping term instead of thermal diffusion, this paper aims to
provide a new mathematical result in Besov spaces. Specifically, it shows that the temperature
enhances dissipation and contributes to stabilizing the fluid. A previous study by Lai et al.\
\cite{MR4790886} demonstrated the global existence and stability of the system \eqref{1.2} in the
Sobolev space $H^3$. Furthermore, they also established large-time behavior for the solutions:
\begin{equation} \label{1.1.1}
\begin{aligned}
 \|(u_2,\theta)\|_{L^2}\to 0,\quad\text{as } t\to\infty, \\
 \|(\nabla u,\nabla b,\nabla\theta)\|_{H^2}\lesssim(1+t) ^{-1/2}.
\end{aligned}
\end{equation}
However, important questions regarding large-time behavior, such as explicit decay rates for the
solution itself and for its high-order derivatives, remains unresolved.
The primary objective of this paper is to establish optimal decay rates for the solutions of the 2D
magnetic B\'enard problem \eqref{1.2}, which align with the decay rates of the heat semi-group.

The purpose of this paper is to study the asymptotic stability near an equilibrium state.
We note that if a solution $\theta_e(x_2)$ satisfies $\theta_e'(x_2^0)<0$ for some
$x_2^0\in\mathbb{R}$, which implies that fluid with a higher temperature lies below that with a lower temperature, then the system is unstable-the Rayleigh--B\'enard instability occurs, as discussed
in \cite{MR128226,MR2098531,MR4134456}, among other. Therefore, we focus on the opposite case, where $\theta_e'(x_2^0)>0$, implying that the fluid with lower temperature lies below the fluid with
higher temperature. More specifically, we consider the equilibrium state
$$
{B}_e=(0,1)^T,\quad \vartheta_e=x_2,\quad
 P_e=\frac{1}{2}x_2^2.
 $$
 The perturbations around this equilibrium state are defined by
\begin{equation*}
 b=B-{B}_e,\quad \theta=\vartheta-\vartheta_e,\quad p=P-P_e.
\end{equation*}
The evolution of the perturbations $(u,b,\theta)$ is governed by the following inviscid system of
equations:
\begin{equation}\label{1.2}
\begin{gathered}
 u_t+ u\cdot \nabla u+\nabla p= (0,\theta)^\mathrm{T}+ b\cdot \nabla b+\partial_2b,\\
 b_t+u\cdot\nabla b-\mu\Delta b=b\cdot\nabla u+\partial_2u,\\
 \theta_t+u\cdot \nabla \theta+\nu\theta=-u_2,\\
 \operatorname{div} u=\operatorname{div} b=0.
\end{gathered}
\end{equation}
 The corresponding initial conditions are
 \begin{equation} \label{1.1.2}
 u(x,0)=u_0(x),\quad b(x,0)=b_0(x),\quad \theta(x,0)=\theta_0(x).
 \end{equation}

 Before presenting our main results, we clarify the decomposition into low-frequency and high-frequency
 components. For any $u\in\mathcal{S}'$, we define its low-frequency part $u^{L}$ and high-frequency
 part $u^{H}$ via the Littlewood--Paley decomposition as
\[
 u^L:=\sum_{j\leqslant-1}\dot{\Delta}_ju,\ \ u^H:=\sum_{j\geqslant 0}\dot{\Delta}_ju.
\]
To measure separately the contributions of low and high frequencies in Besov norms, we use the following notation:
\begin{gather*}
\|u\|_{\dot{B}^{s}_{p,r}}^L:=\|(2^{js}\|\dot{\Delta}_ju\|_{L^p})_{j\leqslant 0}\|_{l^r},\ \ \|u\|_{\dot{B}^{s}_{p,r}}^H:=\|(2^{js}\|\dot{\Delta}_ju\|_{L^p})_{j\geqslant -1}\|_{l^r},\\
\|u\|_{\tilde{L}^q_t(\dot{B}^s_{p,r})}^L:=\|(2^{js}\|\dot{\Delta}_ju\|_{L^q_t(L^p)})_{j\leqslant 0}\|_{l^r},\ \ \|u\|_{\tilde{L}^q_t(\dot{B}^s_{p,r})}^H:=\|(2^{js}\|\dot{\Delta}_ju\|_{L^q_t(L^p)})_{j\geqslant -1}\|_{l^r}.
\end{gather*}
For additional background on Besov spaces and related notation, we refer the reader to Section 2.
Now we state the nonlinear stability result of the magnetic B\'enard problem.


\begin{theorem} \label{thm1}
Let $(u_0^L, b_0^L, \theta_0^L)\in \dot{B}^0_{2,1}(\mathbb{R}^2)$ and
$(u_0^H, b_0^H, \theta_0^H)\in \dot{B}^2_{2,1}(\mathbb{R}^2)$, with
$\operatorname{div} u_0=\operatorname{div} b_0=0$. Then, there exists a
 constant $\epsilon_0>0$ such that if the initial data satisfies
\begin{equation} \label{1.4}
 \mathcal{E}_0:=\|(u_0,b_0,\theta_0)\|_{\dot{B}^0_{2,1}}^L+\|(u_0,b_0,\theta_0)\|_{\dot{B}^2_{2,1}}^H\leqslant\epsilon_0,
\end{equation}
 system \eqref{1.2} and \eqref{1.1.2} admit a unique global solution with
\begin{equation} \label{1.5}
\begin{gathered}
(u^L,b^L,\theta^L)\in C(\mathbb{R}_+;\dot{B}^0_{2,1})\cap L^1(\mathbb{R}_+;\dot{B}^2_{2,1}),\\
 b^H\in C(\mathbb{R}_+;\dot{B}^2_{2,1})\cap L^1(\mathbb{R}_+;\dot{B}^3_{2,1}), \\
(u^H,\theta^H)\in C(\mathbb{R}_+;\dot{B}^2_{2,1})\cap L^1(\mathbb{R}_+;\dot{B}^2_{2,1}).
\end{gathered}
\end{equation}
Moreover, for all $t > 0$, the solution $(u, b, \theta)$ satisfies
\begin{equation}\label{1.3}
\begin{aligned}
&\|(u,b,\theta)\|_{\tilde{L}^\infty_t(\dot{B}^0_{2,1})}^L
 +\|(u,b,\theta)\|_{\tilde{L}^\infty_t(\dot{B}^2_{2,1})}^H
  +\|(u,b,\theta)\|_{{L}^1_t(\dot{B}^2_{2,1})}^L
 +\|(u,\theta)\|_{{L}^1_t(\dot{B}^2_{2,1})}^H+\|b\|_{{L}^1_t(\dot{B}^3_{2,1})}^H\\
&\lesssim\mathcal{E}_0.
\end{aligned}
\end{equation}
\end{theorem}


\begin{remark} \rm
If the fluid is unaffected by temperature and magnetic field, the system \eqref{1.2} reduces to the
2D incompressible Euler equation, where the vorticity gradient can grow double exponentially over
time \cite{MR3293735,MR3245016,MR3475939}.
However, when considering the interaction between temperature and the magnetic field, the system \eqref{1.2} becomes stable.
\end{remark}


\begin{remark} \rm
Theorem \ref{thm1} shows that if the initial data is sufficiently small in $\dot{B}^0_{2,1}$ and
$\dot{B}^2_{2,1}$, then system \eqref{1.2} admits a global solution.
Since $\dot{H}^s(\mathbb{R}^2)\hookrightarrow\dot{B}^2_{2,1}(\mathbb{R}^2)$ for $s>2$,
we can replace $\dot{B}^0_{2,1}$ and $\dot{B}^2_{2,1}$ with $H^3$.
This existence result can be found in \cite{MR4790886}, and here we extend it to the homogeneous
Besov spaces.
\end{remark}

Then we also obtain the optimal temporal decay rates of the solution obtained in Theorem \ref{thm1}.

\begin{theorem}\label{thm2}
Assume that the initial data $(u_0^L,b_0^L,\theta_0^L)\in \dot{B}^{-\sigma}_{2,\infty}(\mathbb{R}^2)(0<\sigma\leqslant 1)$. Then,
the corresponding solution $(u,b,\theta)$ obtained in Theorem \ref{thm1} satisfies the following optimal temporal decay rates for all $t\geqslant 0$:
\begin{equation} \label{1.6}
\|\Lambda^s(u,b,\theta)\|_{L^2}\lesssim(1+t)^{-\frac{\sigma+s}{2}},
\end{equation}
where the pseudo-differential operator 
$\Lambda:=(-\Delta)^{1/2}$ and $s\in(-\sigma,0]$.
\end{theorem}

\begin{remark} \rm
The temporal decay rates given in \eqref{1.6} are optimal in the sense that they 
coincide with that of the heat semi-group. In contrast to the large time behavior 
$\|(u_2,\theta)(t)\|_{L^2}\to 0$  as $t\to\infty$ described in \cite{MR4790886}, 
our result provides a precise decay rate.
 \end{remark}

 If the damping term $\nu\theta$ in \eqref{1.2} is replaced by the thermal diffusivity term
 $\kappa\Delta\theta$, then \eqref{1.2} can be written as
 \begin{equation}\label{1.12}
 \begin{gathered}
 u_t+ u\cdot \nabla u+\nabla p= (0,\theta)^\mathrm{T}+ b\cdot \nabla b+\partial_2b,\\
 b_t+u\cdot\nabla b-\mu\Delta b=b\cdot\nabla u+\partial_2u,\\
 \theta_t+u\cdot \nabla \theta-\kappa\Delta\theta=-u_2,\\
 \operatorname{div} u=\operatorname{div} b=0,\\
 u(x,0)=u_0(x),\ \ b(x,0)=b_0(x),\ \ \theta(x,0)=\theta_0(x).
 \end{gathered}
 \end{equation}

 For the system in \eqref{1.12}, the stability result and optimal temporal decay rates still hold.
 The proof of Theorem \ref{thm3} is similar to that of Theorem \ref{thm1} and \ref{thm2}, so we omit it here.

 \begin{theorem}\label{thm3}
 Let $(u_0^L, b_0^L, \theta_0^L)\in \dot{B}^0_{2,1}(\mathbb{R}^2)$ and
 $(u_0^H, b_0^H, \theta_0^H)\in \dot{B}^2_{2,1}(\mathbb{R}^2)$, with
 $\operatorname{div} u_0=\operatorname{div} b_0=0$. Then, there exists a
constant $\epsilon_0>0$ such that if the initial data satisfies
\begin{equation} \label{1.10}
 \mathcal{E}_0:=\|(u_0,b_0,\theta_0)\|_{\dot{B}^0_{2,1}}^L+\|(u_0,b_0,\theta_0)\|_{\dot{B}^2_{2,1}}^H\leqslant\epsilon_0,
\end{equation}
then system \eqref{1.12} admits a unique global solution with
\begin{equation} \label{1.11}
\begin{gathered}
(u^L,b^L,\theta^L)\in C(\mathbb{R}_+;\dot{B}^0_{2,1})\cap L^1(\mathbb{R}_+;\dot{B}^2_{2,1}),\\
u^H\in C(\mathbb{R}_+;\dot{B}^2_{2,1})\cap L^1(\mathbb{R}_+;\dot{B}^2_{2,1}), \\
(b^H,\theta^H)\in C(\mathbb{R}_+;\dot{B}^2_{2,1})\cap L^1(\mathbb{R}_+;\dot{B}^3_{2,1}).
 \end{gathered}
\end{equation}
 In addition, if the initial data $(u_0^L,b_0^L,\theta_0^L)\in \dot{B}^{-\sigma}_{2,\infty}(\mathbb{R}^2)$
 with $0<\sigma\leqslant 1$, then
 the following optimal temporal decay rates hold for all $t\geqslant 0$:
 \begin{equation} \label{1.16}
 \|\Lambda^s(u,b,\theta)\|_{L^2}\lesssim(1+t)^{-\frac{\sigma+s}{2}}.
 \end{equation}
 \end{theorem}

Now we outline the main idea of the proof. The key to proving the global existence of the solution is to
establish uniform \emph{a priori} estimates. A significant challenge arises in controlling the convection
term $u\cdot\nabla u$ due to the absence of viscosity in the system. To overcome this, we exploit the
potential hyperbolic structure of the system. By applying the Leray--Helmholtz projection
$\mathbb{P}:=I-\nabla\Delta^{-1}{\rm div}$ to the first equation of \eqref{1.2}, we obtain the system
\begin{equation} \label{1.1.3}
\begin{gathered}
\partial_tu_1+\partial_1\partial_2\Delta^{-1}\theta-\partial_2b_1=\mathbb{P}(b\cdot\nabla b_1-u\cdot\nabla u_1),\\
\partial_tu_2-\partial_1^2\Delta^{-1}\theta-\partial_2b_2=\mathbb{P}(b\cdot\nabla b_2-u\cdot\nabla u_2),
\end{gathered}
\end{equation}
where $I$ represents the identity matrix.
If the temperature is neglected, differentiating \eqref{1.2} with respect to time and making several substitutions, using \eqref{1.1.3}, we have
\begin{equation} \label{1.7}
\begin{gathered}
\partial_{tt}u_1-\mu\partial_t\Delta u_1-\partial_2^2u_1=\mathcal{M}_1,\\
\partial_{tt}u_2-\mu\partial_t\Delta u_2-\partial_2^2u_2=\mathcal{M}_2,\\
\partial_{tt}b-\mu\partial_t\Delta b-\partial_2^2b=\mathcal{M}_3,\\
\end{gathered}
\end{equation}
where
\begin{gather*}
\mathcal{M}_1:=(\partial_t-\mu\Delta)\mathbb{P}(b\cdot\nabla b_1-u\cdot\nabla u_1)+\partial_2(b\cdot\nabla u_1-u\cdot\nabla b_1),\\
\mathcal{M}_2:=(\partial_t-\mu\Delta)\mathbb{P}(b\cdot\nabla b_2-u\cdot\nabla u_2+\theta)+\partial_2(b\cdot\nabla u_2-u\cdot\nabla b_2),\\
\mathcal{M}_3:=\partial_t(b\cdot\nabla u-u\cdot\nabla b)+\partial_2\mathbb{P}(b\cdot\nabla b-u\cdot\nabla u+(0,\theta)^\mathrm{T}).
\end{gather*}
Leveraging the hyperbolic smoothing effect of the above equations, we observe that the magnetic
field captures a vertical dissipation effect on the velocity (see \eqref{3.8} for details), that is,
$\partial_2u$. Similarly, if the fluid is not affected by the magnetic field, differentiating \eqref{1.2} with respect to time and making several substitutions, we also obtain
\begin{equation} \label{1.1.4}
\begin{gathered}
 \partial_{tt}u_1+\nu\partial_tu_1+\partial_1^2\Delta^{-1}u_1=\mathcal{M}_4,\\
 \partial_{tt}u_2+\nu\partial_t u_2+\partial_1^2\Delta^{-1}u_2=\mathcal{M}_5,\\
 \partial_{tt}\theta+\nu\partial_t\theta+\partial_1^2\Delta^{-1}\theta=\mathcal{M}_6,\\
 \end{gathered}
\end{equation}
 where
\begin{gather*}
\mathcal{M}_4:=(\partial_t+\nu)\mathbb{P}(b\cdot\nabla b_1-u\cdot\nabla u_1+\partial_2b_1)+\partial_1\partial_2\Delta^{-1}(u\cdot\nabla\theta),\\
\mathcal{M}_5:=(\partial_t-\nu)\mathbb{P}(b\cdot\nabla b_2-u\cdot\nabla u_2+\partial_2b_2)-\partial_1^2\Delta^{-1}(u\cdot\nabla\theta),\\
\mathcal{M}_6:=\mathbb{P}(u\cdot\nabla u_2-b\cdot\nabla b_2)+\partial_t(u\cdot\nabla \theta)-\partial_2b_2.
\end{gather*}
This system also suggests a potential construction of a hyperbolic structure. It shows that the
temperature contributes a horizontal dissipation effect on the velocity (see \eqref{3.6} for details),
that is, $\partial_1u$. Based on the above observations, to handle the nonlinear terms, we employ
the technique of high and low frequency decomposition, separating the solution $(u,b,\theta)$ into
its low-frequency and high-frequency parts. For the low-frequency part, we focus on ensuring the
minimal regularity requirement, i.e., $(u^L,b^L,\theta^L)\in\dot{B}^0_{2,1}$. For the high-frequency part,
we control the nonlinear interactions using appropriate product and commutator estimates.
This leads to the functional framework $(u^H,b^H,\theta^H)\in\dot{B}^2_{2,1}$. Finally, by applying interpolation inequality and negative Besov norms, we derive a Lyapunov-type differential inequality,
which leads to the optimal temporal decay rates.

This article paper is organized as follows. In Section 2, we introduce the Littlewood-Paley
decomposition, Besov spaces and other essential tools. In Section 3, we prove
Theorem \ref{thm1}, and in Section 4, we establish the proof of Theorem \ref{thm2}.


\subsection*{Notation} For $1< p\leqslant+\infty$, we simplify $\int_{\mathbb{R}^2}$,
$\|f\|_{L^p(\mathbb{R}^2)}$ and $\dot{B}^p_{q,r}(\mathbb{R}^2)$ as
$\int$, $\|f\|_{L^p}$ and $\dot{B}^p_{q,r}$, respectively. For a Banach space $A$, we use the
shorthand notation $\|(f,g)\|_A:=\|f\|_A+\|g\|_A$.
 The symbols $C$ and $C_i$ $(i=1,2,3,4)$ represent generic positive constants, which may vary
 from one line to another. For a uniform constant $C$, we write $f\lesssim g$ to mean $f\leqslant Cg$,
 and $f\approx g$ to indicate both $f\leqslant Cg$ and $g\leqslant Cf$. For two operators $X$ and $Y$, we define the commutator as $[X,Y]=XY-YX$. $\mathcal{F}$ ($\mathcal{F}^{-1}$) is the Fourier (inverse Fourier) transform operator.

\section{Preliminaries}\label{subsec:02}

In this section, we introduce the Littlewood-Paley decomposition, the definition of
Besov spaces and some useful properties. For further details, we refer to
\cite{MR2768550}.

Let $\chi$ be radial function such that $\chi\equiv1$ in
$\{\xi\in\mathbb{R}^2:|\xi|\leqslant 3/4\}$ which supported in
$\{\xi\in\mathbb{R}^2:|\xi|\leqslant 1\}$. Then $\varphi:=\chi(\xi)-\chi(2\xi)$
is supported in the annulus $\{\xi\in\mathbb{R}^2:3/4\leqslant|\xi|\leqslant 8/3\}$
and satisfies
\begin{gather*}
 \chi(\xi)+\sum_{j\geqslant 0}\varphi(2^{-j}\xi)=1,\quad \forall\xi\in\mathbb{R}^2,\\
 \sum_{j\in\mathbb{Z}}\varphi(2^{-j}\xi)=1,\quad \forall\xi\in\mathbb{R}^2\setminus \{0\}.
\end{gather*}
The homogeneous dyadic blocks $\dot{\Delta}_j$ are defined by
\[
 \dot{\Delta}_ju:=\varphi(2^{-j}D)u=2^{2j}\int h(2^jy)u(x-y){\rm d}y,\quad
 \forall j\in\mathbb{Z},
\]
where $h:=\mathcal{F}^{-1}\varphi$. For tempered distribution $u\in \mathcal{S}'$,
 we have the decomposition
\[
 u=\sum_{j\in\mathbb{Z}}\dot{\Delta}_ju.
\]
According to the above decomposition, we have the definition of the homogeneous Besov
spaces.
\begin{definition}\label{def2.1} \rm
 For $s\in\mathbb{R}$ and $1\leqslant p,r\leqslant\infty$, then the homogeneous
Besov space $\dot{B}^s_{p,r}$ is defined as
\[
 \dot{B}^s_{p,r}:=\{u\in\mathcal{S}':\|u\|_{\dot{B}^s_{p,r}}<\infty\},
\]
 where
\[
 \|u\|_{\dot{B}^s_{p,r}}:=\|(2^{js}\|\dot{\Delta}_ju\|_{L^p})_{j\in\mathbb{Z}}\|_{l^r}.
\]
\end{definition}

Next, we present Berstein's inequalities, which will be used frequently.

\begin{lemma} \label{lem2.1}
For any $r\in(0,R)$, nonnegative integer $k$ and pair $(p,q)\in[1,\infty]^2$ with
$p\in[1,q]$, there exists a constant $C>0$ for $u\in L^p$ such that
\begin{gather*}
 \operatorname{supp}\mathcal{F}u\subset \{\xi\in\mathbb{R}^2:
 |\xi|\leqslant\lambda R\}
 \Rightarrow\sup _{|\alpha|=k}\|\partial^\alpha u\|_{L^q}
 \leqslant C^{k+1}\lambda^{k+2(\frac{1}{p}-\frac{1}{q})}\|u\|_{L^p},\\
 \operatorname{supp}\mathcal{F}u\subset\{\xi\in\mathbb{R}^2:
 \lambda r\leqslant|\xi|\leqslant\lambda R\}\Rightarrow C^{-k-1}
 \lambda^k\|u\|_{L^p}\leqslant\sup _{|\alpha|=k}
 \|\partial^\alpha u\|_{L^p}\leqslant C^{k+1}\lambda^{k}\|u\|_{L^p}.
 \end{gather*}
\end{lemma}


\begin{lemma}\label{lem2.2}
Assume $s\in \mathbb{R}$ and $1\leqslant p,r\leqslant\infty$, then we have the
following properties:
 \begin{itemize}
\item[(1)] Embedding: for any $1\leqslant r \leqslant\tilde{r}\leqslant\infty$ and
$1\leqslant p \leqslant\tilde{p}\leqslant\infty$, it holds
\begin{equation} \label{2.1}
\dot{B}^s_{p,r}(\mathbb{R}^2)\hookrightarrow\dot{B}^{s-2(\frac{1}{p}
-\frac{1}{\tilde{p}})}_{\tilde{p},\tilde{r}}(\mathbb{R}^2),\quad
\dot{B}^0_{p,1}(\mathbb{R}^2)\hookrightarrow L^p(\mathbb{R}^2), \quad
\dot{B}^{\frac{2}{p}}_{p,1}(\mathbb{R}^2)\hookrightarrow L^\infty(\mathbb{R}^2).
\end{equation}

\item[(2)] Interpolation: for any $0<\theta<1$ and $s_1<s_2$, then it holds that
\begin{gather}
\|u\|_{\dot{B}^{ s_1\theta+s_2(1-\theta)}_{p,r}}
 \leqslant \|u\|_{\dot{B}^{s_1}_{p,r}}^\theta\|u\|_{\dot{B}^{s_2}_{p,r}}^{1-\theta}, \label{2.2}\\
\|u\|_{\dot{B}^{ s_1\theta+s_2(1-\theta)}_{p,1}}
 \leqslant\frac{C}{s_2-s_1}\Big(\frac{1}{\theta}+\frac{1}{1-\theta}\Big)
 \|u\|_{\dot{B}^{s_1}_{p,\infty}}^\theta\|u\|_{\dot{B}^{s_2}_{p,\infty}}^{1-\theta}.
 \label{2.3}
\end{gather}

\item[(3)] Derivatives: for any integer $k$, then it holds that
\begin{equation} \label{2.4}
 \sup _{|\alpha|=k}\|\partial^\alpha u\|_{\dot{B}^{s}_{p,r}}
 \approx\|u\|_{\dot{B}^{s+k}_{p,r}}.
\end{equation}
 \end{itemize}
\end{lemma}

Next, for the sake of time integration, we provide the following mixed Besov space
about time and space, which was introduced in \cite{MR1354312}.

\begin{definition}\label{def2.2} \rm
Assume $s\in\mathbb{R}$, $t>0$ and $1\leqslant p,q,r\leqslant\infty$, the space
$\tilde{L}^q(0,t;\dot{B}^s_{p,r})$ is defined as
\[
 \tilde{L}^q(0,t;\dot{B}^s_{p,r}):=\{u\in L^q(0,t;\mathcal{S}'(\mathbb{R}^2)):\|u\|_{\tilde{L}^q_t(\dot{B}^s_{p,r})}<\infty\},
\]
 where
\[
 \|u\|_{\tilde{L}^q_t(\dot{B}^s_{p,r})}:=\|(2^{js}\|\dot{\Delta}_ju\|_{L^q_t(L^p)})_{j\in\mathbb{Z}}\|_{l^r}.
\]
\end{definition}

\begin{remark}\label{rem2.1} \rm
One easily from Minkowski's inequality check that
 \begin{equation} \label{2.5}
 \begin{gathered}
 \|u\|_{{L}^q_t(\dot{B}^s_{p,r})}\leqslant\|u\|_{\tilde{L}^q_t(\dot{B}^s_{p,r})},
 \quad\text{if }r\leqslant q,\\
 \|u\|_{\tilde{L}^q_t(\dot{B}^s_{p,r})}\leqslant\|u\|_{{L}^q_t(\dot{B}^s_{p,r})},
 \quad\text{if }q\leqslant r.
 \end{gathered}
 \end{equation}
\end{remark}

For any $s'>0$, one observes that
\begin{equation} \label{2.6}
 \begin{gathered}
 \|u^L\|_{\dot{B}^s_{p,r}}\lesssim\|u\|^L_{\dot{B}^s_{p,r}}\lesssim\|u\|^L_{\dot{B}^{s-s'}_{p,r}},\\
 \|u^H\|_{\dot{B}^s_{p,r}}\lesssim\|u\|^H_{\dot{B}^s_{p,r}}\lesssim\|u\|^H_{\dot{B}^{s+s'}_{p,r}}.
 \end{gathered}
\end{equation}

Finally, let us recall some product estimates and commutator estimates.

\begin{lemma}\label{lem2.3}
 We have the following product estimates:
 \begin{itemize}
\item For any $1\leqslant p,r\leqslant\infty$ and $s>0$, it holds that
 \begin{equation}\label{2.7}
 \|fg\|_{\dot{B}^s_{p,r}}\lesssim\|f\|_{\dot{B}^s_{p,r}}\|g\|_{L^\infty}+\|g\|_{\dot{B}^s_{p,r}}\|f\|_{L^\infty}.
 \end{equation}

\item For any $s_1\leqslant\frac{2}{p}$, $s_2\leqslant\frac{2}{p}$, $s_1+s_2>0$ and
$1\leqslant p\leqslant\infty$, it holds that
\begin{equation} \label{2.8}
 \|fg\|_{\dot{B}^{s_1+s_2-\frac{2}{p}}_{p,1}}\lesssim\|f\|_{\dot{B}^{s_1}_{p,1}}\|g\|_{\dot{B}^{s_2}_{p,1}}.
\end{equation}

\item For any $s_1\leqslant\frac{2}{p}$, $s_2<\frac{2}{p}$, $s_1+s_2\geqslant0$ and $2\leqslant p\leqslant\infty$, it holds that
\begin{equation}\label{2.9}
 \|fg\|_{\dot{B}^{s_1+s_2-\frac{2}{p}}_{p,\infty}}\lesssim\|f\|_{\dot{B}^{s_1}_{p,1}}\|g\|_{\dot{B}^{s_2}_{p,\infty}}.
\end{equation}
 \end{itemize}
\end{lemma}

\begin{lemma}\label{lem2.4}
For $1\leqslant p\leqslant\infty$ and a vector field $X:=(X_1,X_2)^T$, we have the
following commutator estimates:
 \begin{itemize}
\item For $-\frac{2}{p}<s\leqslant1+\frac{2}{p}$, one has
 \begin{equation}\label{2.10}
 \sum_{j\in\mathbb{Z}}2^{js}\|[\dot{\Delta}_j,X\cdot\nabla]f\|_{L^p}\lesssim\|f\|_{\dot{B}^s_{p,1}}\|X\|_{\dot{B}^{1+\frac{2}{p}}_{p,1}}.
 \end{equation}

\item For $s>0$, one has
 \begin{equation}\label{2.11}
 \sum_{j\in\mathbb{Z}}2^{js}\|[\dot{\Delta}_j,X\cdot\nabla]f\|_{L^p}\lesssim\|f\|_{\dot{B}^s_{p,1}}\|\nabla X\|_{L^\infty}+\|X\|_{\dot{B}^{s}_{p,1}}\|\nabla f\|_{L^\infty}.
 \end{equation}

\item For $-\frac{2}{p}\leqslant s\leqslant1+\frac{2}{p}$, one has
 \begin{equation}\label{2.12}
 \sum_{j\in\mathbb{Z}}2^{js}\|[\dot{\Delta}_j,X\cdot\nabla]f\|_{L^p}\lesssim\|f\|_{\dot{B}^s_{p,\infty}}\|X\|_{\dot{B}^{1+\frac{2}{p}}_{p,1}}.
 \end{equation}
\end{itemize}
\end{lemma}

\section{Proof of Theorem \ref{thm1}}

 In this section, we prove the existence of global solutions to system
 \eqref{1.2}, that is, Theorem \ref{thm1}. Combining the local well-posedness with
some a priori estimates, and then employing a standard continuity argument to extend
the local solution to a global one. The proof of local existence for equations
\eqref{1.2} is standard, and we omit the details here. Interested reads are referred
to \cite{MR1855277,MR4496608}. With the local existence result in hand, the key step
is to prove uniform a priori estimates of the solution, as stated in \eqref{1.3}.
To do so, we introduce the energy functionals
 \begin{equation}\label{3.1}
\begin{gathered}
 \mathcal{E}(t)=\|(u,b,\theta)\|^L_{\dot{B}^0_{2,1}}
 +\|(u,b,\theta)\|^H_{\dot{B}^2_{2,1}}, \\
 \mathcal{D}(t)=\|(u,b,\theta)\|^L_{\dot{B}^2_{2,1}}+\|(u,\theta)\|^H_{\dot{B}^2_{2,1}}
 +\|b\|^H_{\dot{B}^3_{2,1}}.
\end{gathered}
\end{equation}
First, we focus on the uniform stability estimate for the low-frequency part of
$(u,b,\theta)$.

\begin{lemma}\label{lem3.1}
It holds that
\begin{equation}\label{3.2}
 \|(u,b,\theta)\|_{\tilde{L}^\infty_t(\dot{B}^0_{2,1})}^L+\|(u,b,\theta)\|_{{L}^1_t(\dot{B}^2_{2,1})}^L\lesssim\|(u_0,b_0,\theta_0)\|_{\dot{B}^0_{2,1}}^L+\int_0^t\mathcal{E}(s)\mathcal{D}(s){\rm d}s.
\end{equation}
\end{lemma}

\begin{proof}
Applying the operator $\dot{\Delta}_j (j\in\mathbb{Z})$ to \eqref{1.2} yields
\begin{equation}\label{3.3}
\begin{gathered}
 \partial_t\dot{\Delta}_ju+\dot{\Delta}_j\nabla p-\dot{\Delta}_j (0,\theta)^\mathrm{T}-\dot{\Delta}_j\partial_2b=\dot{\Delta}_j \mathcal{N}_1,\\
 \partial_t\dot{\Delta}_jb-\mu\Delta \dot{\Delta}_jb-\dot{\Delta}_j\partial_2u=\dot{\Delta}_j\mathcal{N}_2,\\
 \partial_t\dot{\Delta}_j\theta+\nu\dot{\Delta}_j\theta+\dot{\Delta}_ju_2=-\dot{\Delta}_j \mathcal{N}_3,
\end{gathered}
\end{equation}
where
\[
 \mathcal{N}_1:=b\cdot \nabla b-u\cdot \nabla u, \quad
 \mathcal{N}_2:=b\cdot \nabla u-u\cdot \nabla b, \quad
 \mathcal{N}_3:=u\cdot \nabla \theta.
\]

Taking the $L^2$ inner product with $\dot{\Delta}_ju$, $\dot{\Delta}_jb$ and
$\dot{\Delta}_j\theta$ to the first three equations of \eqref{3.3}, respectively,
then using integration by parts and summing up the resulting equality, we obtain
\begin{equation}\label{3.4}
\begin{aligned}
 &\frac{1}{2}\frac{\rm d}{{\rm d}t}\|(\dot{\Delta}_ju,\dot{\Delta}_jb,
 \dot{\Delta}_j\theta)\|_{L^2}^2+\mu\|\nabla\dot{\Delta}_jb\|_{L^2}^2
 +\nu\|\dot{\Delta}_j\theta\|_{L^2}^2 \\
 &=\int\Big(\dot{\Delta}_j \mathcal{N}_1\cdot \dot{\Delta}_ju+\dot{\Delta}_j
 \mathcal{N}_2\cdot \dot{\Delta}_jb-\dot{\Delta}_j \mathcal{N}_3
 \dot{\Delta}_j\theta\Big){\rm d}x,
\end{aligned}
\end{equation}
where we have used the cancelation that
\[
 \int\dot{\Delta}_j(\nabla p-(0,\theta)-\partial_2b)\cdot\dot{\Delta}_j u{\rm d}x-\int\dot{\Delta}_j\partial_2u\cdot\dot{\Delta}_jb{\rm d}x+\int\dot{\Delta}_ju_2\dot{\Delta}_j\theta{\rm d}x=0.
\]

So as to get the horizontal dissipation of $u$, applying the operator
$\dot{\Delta}_j\nabla$ to the second component of the momentum equation $\eqref{1.2}_1$
and $\eqref{1.2}_3$ lead to
\begin{equation}\label{3.5}
\begin{gathered}
 \partial_t\dot{\Delta}_j\nabla u_2+\dot{\Delta}_j\nabla\partial_2 p-\dot{\Delta}_j\nabla \theta-\dot{\Delta}_j\nabla\partial_2b_2=\dot{\Delta}_j\nabla \mathcal{N}_1^2,\\
 \partial_t\dot{\Delta}_j\nabla\theta+\nu\dot{\Delta}_j\nabla\theta+\dot{\Delta}_j\nabla u_2=-\dot{\Delta}_j\nabla\mathcal{N}_3,
\end{gathered}
\end{equation}
where $\mathcal{N}_1^2$ stands for the second component of $\mathcal{N}_1$.

Noting that ${\rm div}u=0$, then we have
$\|\dot{\Delta}_j\partial_1u\|_{L^2}=\|\dot{\Delta}_j\nabla u_2\|_{L^2}$.
Thus, multiplying $\eqref{3.5}_1$ and $\eqref{3.5}_2$ by
$\dot{\Delta}_j\nabla\theta$ and $\dot{\Delta}_j\nabla u_2$ in $L^2$, respectively,
we find that
\begin{equation}\label{3.6}
\begin{aligned}
&\frac{\rm d}{{\rm d}t}\int\dot{\Delta}_j\nabla u_2\cdot\dot{\Delta}_j\nabla\theta
{\rm d}x +\|\dot{\Delta}_j\partial_1 u\|_{L^2}-\|\dot{\Delta}_j\nabla \theta\|_{L^2}\\
&+\int\Big(\dot{\Delta}_j\nabla\partial_2p-\dot{\Delta}_j\nabla\partial_2b_2\Big)
 \cdot\dot{\Delta}_j\nabla\theta{\rm d}x
 +\nu\int\dot{\Delta}_j\nabla\theta\cdot\dot{\Delta}_j\nabla u_2{\rm d}x \\
&=\int\Big(\dot{\Delta}_j\nabla \mathcal{N}_1^2\cdot\dot{\Delta}_j\nabla\theta
 -\dot{\Delta}_j\nabla\mathcal{N}_3\cdot\dot{\Delta}_j\nabla u_2\Big){\rm d}x.
\end{aligned}
\end{equation}
To control the term $\int\dot{\Delta}_j\nabla\partial_2p\cdot\dot{\Delta}_j
 \nabla\theta{\rm d}x$, applying the operator ${\rm div}$ to $\eqref{1.2}_1$, by
$\eqref{1.2}_4$ we check that
\begin{equation}\label{3.7}
 \Delta p=\nabla b:\nabla b-\nabla u:\nabla u+\partial_2\theta.
\end{equation}

Next, we aim to obtain the vertical dissipation of $u$. Taking $L^2$ inner product of
$\eqref{3.3}_1$ and $\eqref{3.3}_2$ with $\dot{\Delta}_j\partial_2b$ and
$-\dot{\Delta}_j\partial_2u$, respectively, we obtain
\begin{equation}\label{3.8}
\begin{aligned}
&\frac{\rm d}{{\rm d}t}\int\dot{\Delta}_j u\cdot\dot{\Delta}_j\partial_2b{\rm d}x
 +\|\dot{\Delta}_j\partial_2 u\|_{L^2}-\|\dot{\Delta}_j\partial_2b\|_{L^2}
 -\int\dot{\Delta}_j\theta\dot{\Delta}_j\partial_2b_2{\rm d}x
 +\nu\int\Delta\dot{\Delta}_jb\cdot\dot{\Delta}_j\partial_2 u{\rm d}x \\
&=\int\Big(\dot{\Delta}_j \mathcal{N}_1\cdot\dot{\Delta}_j\partial_2b
 -\dot{\Delta}_j\mathcal{N}_2\cdot\dot{\Delta}_j\partial_2 u\Big){\rm d}x,
\end{aligned}
\end{equation}
where we have used $\int\dot{\Delta}_j\nabla p\cdot\dot{\Delta}_j\partial_2b{\rm d}x=0$.

Then combining \eqref{3.4}, \eqref{3.6}, \eqref{3.7} and \eqref{3.8}, for a large
constant $C_1$ we can obtain the estimate
\begin{equation}\label{3.9}
\begin{aligned}
\frac{\rm d}{{\rm d}t}\tilde{\mathcal{E}}_L^2+\tilde{\mathcal{D}}_L^2
&=C_1\int\dot{\Delta}_j \mathcal{N}_1\cdot \dot{\Delta}_ju{\rm d}x+C_1\int\dot{\Delta}_j \mathcal{N}_2\cdot \dot{\Delta}_jb{\rm d}x-C_1\int\dot{\Delta}_j \mathcal{N}_3\cdot \dot{\Delta}_j\theta{\rm d}x \\
&\quad +\int\dot{\Delta}_j\nabla \mathcal{N}_1^2\cdot\dot{\Delta}_j\nabla\theta{\rm d}x-\int\dot{\Delta}_j\nabla\mathcal{N}_3\cdot\dot{\Delta}_j\nabla u_2{\rm d}x+\int\dot{\Delta}_j \mathcal{N}_1\cdot\dot{\Delta}_j\partial_2b{\rm d}x \\
&\quad -\int\dot{\Delta}_j\mathcal{N}_2\cdot\dot{\Delta}_j\partial_2 u{\rm d}x+\int\dot{\Delta}_j(\nabla u:\nabla u-\nabla b:\nabla b)\dot{\Delta}_j\partial_2\theta{\rm d}x,
\end{aligned}
\end{equation}
 where
 \begin{gather*}
 \tilde{\mathcal{E}}_L^2:=\frac{1}{2}C_1\|(\dot{\Delta}_ju,\dot{\Delta}_jb,
 \dot{\Delta}_j\theta)\|_{L^2}^2+\int\dot{\Delta}_j\nabla u_2\cdot\dot{\Delta}_j
 \nabla\theta{\rm d}x
 +\int\dot{\Delta}_j u\cdot\dot{\Delta}_j\partial_2b{\rm d}x,\\
\begin{aligned}
\tilde{\mathcal{D}}_L^2
&:= C_1\mu\|\nabla\dot{\Delta}_jb\|_{L^2}^2+C_1\nu\|\dot{\Delta}_j\theta\|_{L^2}^2+\|\dot{\Delta}_j\partial_1 u\|_{L^2}^2+\|\dot{\Delta}_j\partial_2 u\|_{L^2}^2+\|\dot{\Delta}_j\partial_2\theta\|_{L^2}^2 \\
&\quad -\|\dot{\Delta}_j\partial_2b\|_{L^2}^2-\|\dot{\Delta}_j\nabla \theta\|_{L^2}^2
 -\int\dot{\Delta}_j\nabla\partial_2b_2\cdot\dot{\Delta}_j\nabla\theta{\rm d}x
 +\nu\int\dot{\Delta}_j\nabla\theta\cdot\dot{\Delta}_j\nabla u_2{\rm d}x \\
&\quad -\int\dot{\Delta}_j\theta\dot{\Delta}_j\partial_2b_2{\rm d}x
 +\nu\int\Delta\dot{\Delta}_jb\cdot\dot{\Delta}_j\partial_2 u{\rm d}x.
\end{aligned}
\end{gather*}
By using Berstein's inequality, for $j\leqslant 0$, one has
\[
 \|\nabla\dot{\Delta}_jf\|_{L^2}\lesssim2^j\|\dot{\Delta}_jf\|_{L^2}
 \lesssim\|\dot{\Delta}_jf\|_{L^2}.
\]
Then taking $C_1$ large enough and using the above result, we have
 \begin{equation}\label{3.10}
\begin{gathered}
 \tilde{\mathcal{E}}_L^2\approx\|(\dot{\Delta}_ju,\dot{\Delta}_jb,\dot{\Delta}_j\theta)
 \|_{L^2}^2, \\
 \tilde{\mathcal{D}}_L^2\approx\|(\nabla\dot{\Delta}_ju,\nabla\dot{\Delta}_jb,
 \dot{\Delta}_j\theta)\|_{L^2}^2\gtrsim 2^{2j}\|(\dot{\Delta}_ju,\dot{\Delta}_jb,
 \dot{\Delta}_j\theta)\|_{L^2}^2.
\end{gathered}
\end{equation}
Integration by parts,and using  H\"older's inequality and Berstein's inequality,
for $j\leqslant 0$, \eqref{3.9} and \eqref{3.10} we have
\[
 \frac{\rm d}{{\rm d}t}\tilde{\mathcal{E}}_L^2+2^{2j}
 \tilde{\mathcal{E}}_L^2\lesssim\|(\dot{\Delta}_j 
 \mathcal{N}_1,\dot{\Delta}_j \mathcal{N}_2,
 \dot{\Delta}_j \mathcal{N}_3,\dot{\Delta}_j (\nabla u:\nabla u-\nabla b:\nabla b))
 \|_{L^2}\tilde{\mathcal{E}}_L,
\]
which together with $\|\dot{\Delta}_j (\nabla u:\nabla u
-\nabla b:\nabla b))\|_{L^2}\lesssim\|\dot{\Delta}_j\mathcal{N}_1\|_{L^2}$
yields
\begin{equation}\label{3.11}
 \tilde{\mathcal{E}}_L(t)+2^{2j}\int_0^t\tilde{\mathcal{E}}_L(s){\rm d}s\lesssim\tilde{\mathcal{E}}_L(0)+\int_0^t\|(\dot{\Delta}_j \mathcal{N}_1,\dot{\Delta}_j \mathcal{N}_2,\dot{\Delta}_j \mathcal{N}_3)\|_{L^2}{\rm d}s.
\end{equation}
Summing over $j\leqslant 0$, together with \eqref{3.10} one obtains
 \begin{equation}\label{3.12}
\begin{aligned}
&\|(u,b,\theta)\|_{\tilde{L}^\infty_t(\dot{B}^0_{2,1})}^L+\|(u,b,\theta)
\|_{{L}^1_t(\dot{B}^2_{2,1})}^L \\
&\lesssim\|(u_0,b_0,\theta_0)\|_{\dot{B}^0_{2,1}}^L+\int_0^t\|( \mathcal{N}_1,
\mathcal{N}_2, \mathcal{N}_3)\|_{\dot{B}^0_{2,1}}^L{\rm d}s.
\end{aligned}
\end{equation}

Now, we estimate the nonlinear terms in the right-hand side of \eqref{3.12}.
According to \eqref{2.4}, \eqref{2.6} and \eqref{2.8}, we obtain
\begin{equation}\label{3.13}
 \|(u\cdot\nabla u,u\cdot\nabla b,u\cdot\nabla \theta)\|_{\dot{B}^0_{2,1}}^L
 \lesssim \|u\|_{{\dot{B}^0_{2,1}}}\|\nabla(u,b,\theta)\|_{{\dot{B}^1_{2,1}}}
 \lesssim \|u\|_{{\dot{B}^0_{2,1}}}\|(u,b,\theta)\|_{{\dot{B}^2_{2,1}}}
 \lesssim \mathcal{E}(t)\mathcal{D}(t),
\end{equation}
and
\begin{equation}\label{3.131}
 \|(b\cdot\nabla b,b\cdot\nabla u)\|_{\dot{B}^0_{2,1}}^L
 \lesssim \|b\|_{{\dot{B}^0_{2,1}}}\|\nabla(b,u)\|_{{\dot{B}^1_{2,1}}}
 \lesssim \|b\|_{{\dot{B}^0_{2,1}}}\|(b,u)\|_{{\dot{B}^2_{2,1}}}
 \lesssim \mathcal{E}(t)\mathcal{D}(t).
\end{equation}
In summary, putting the nonlinear estimates \eqref{3.13} and \eqref{3.131} into
\eqref{3.12}, we obtain the desired estimate \eqref{3.2}. The proof is complete.
 \end{proof}


Then we begin to control the stability estimate for the high-frequency part of
$(u,b,\theta)$.

\begin{lemma}\label{lem3.2}
It holds
 \begin{equation}\label{3.14}
\begin{aligned}
&\|(u,b,\theta)\|_{\tilde{L}^\infty_t(\dot{B}^2_{2,1})}^H
 +\|(u,\theta)\|_{{L}^1_t(\dot{B}^2_{2,1})}^H+\|b\|_{{L}^1_t(\dot{B}^3_{2,1})}^H \\
&\lesssim\|(u_0,b_0,\theta_0)\|_{\dot{B}^2_{2,1}}^H
 +\int_0^t\mathcal{E}(s)\mathcal{D}(s){\rm d}s.
\end{aligned}
\end{equation}
\end{lemma}

\begin{proof}
Multiplying $\eqref{3.3}_1$, $\eqref{3.3}_2$ and $\eqref{3.3}_3$ by
$ -\Delta\dot{\Delta}_j u$, $ -\Delta\dot{\Delta}_j b$ and
$ -\Delta\dot{\Delta}_j \theta$ in $L^2$, respectively, by integrating by parts, the fact that
\[
 \int\dot{\Delta}_j(\partial_2b+(0,\theta)-\nabla p)\cdot\dot\Delta{\Delta}_j u{\rm d}x
 +\int\dot{\Delta}_j\partial_2u\cdot\Delta\dot{\Delta}_j b{\rm d}x
 -\int\dot{\Delta}_ju_2\Delta\dot{\Delta}_j\theta{\rm d}x=0,
\]
and summing up the resulting equality, we have
\begin{equation}\label{3.15}
\begin{aligned}
 &\frac{1}{2}\frac{\rm d}{{\rm d}t}\|(\nabla\dot{\Delta}_ju,\nabla\dot{\Delta}_jb,\nabla\dot{\Delta}_j\theta)\|_{L^2}^2+\mu\|\Delta\dot{\Delta}_jb\|_{L^2}^2+\nu\|\nabla\dot{\Delta}_j\theta\|_{L^2}^2 \\
 &=\int\Big(\nabla\dot{\Delta}_j \mathcal{N}_1\cdot\nabla \dot{\Delta}_ju+\nabla
  \dot{\Delta}_j \mathcal{N}_2\cdot \nabla\dot{\Delta}_jb-\nabla\dot{\Delta}_j
  \mathcal{N}_3\cdot \nabla\dot{\Delta}_j\theta\Big){\rm d}x.
\end{aligned}
\end{equation}
Multiplying \eqref{3.15} by a large constant $C_2$ and adding to \eqref{3.6}, \eqref{3.7}
and \eqref{3.8}, we obtain
 \begin{equation}\label{3.16}
\begin{aligned}
&\frac{\rm d}{{\rm d}t}\tilde{\mathcal{E}}_H^2+\tilde{\mathcal{D}}_H^2\\
&= C_2\int\nabla\dot{\Delta}_j \mathcal{N}_1\cdot \nabla\dot{\Delta}_ju{\rm d}x
 +C_2\int\nabla\dot{\Delta}_j \mathcal{N}_2\cdot\nabla \dot{\Delta}_jb{\rm d}x
 -C_2\int\nabla\dot{\Delta}_j \mathcal{N}_3\cdot \nabla\dot{\Delta}_j\theta{\rm d}x \\
&\quad +\int\dot{\Delta}_j\nabla \mathcal{N}_1^2\cdot\dot{\Delta}_j\nabla\theta{\rm d}x
 -\int\dot{\Delta}_j\nabla\mathcal{N}_3\cdot\dot{\Delta}_j\nabla u_2{\rm d}x
 +\int\dot{\Delta}_j \mathcal{N}_1\cdot\dot{\Delta}_j\partial_2b{\rm d}x \\
&\quad -\int\dot{\Delta}_j\mathcal{N}_2\cdot\dot{\Delta}_j\partial_2 u{\rm d}x
 +\int\dot{\Delta}_j(\nabla u:\nabla u-\nabla b:\nabla b)
 \dot{\Delta}_j\partial_2\theta{\rm d}x,
\end{aligned}
\end{equation}
where
\begin{gather*}
 \tilde{\mathcal{E}}_H^2
:=\frac{1}{2}C_2\|(\nabla\dot{\Delta}_ju,\nabla\dot{\Delta}_jb,
 \nabla\dot{\Delta}_j\theta)\|_{L^2}^2+\int\dot{\Delta}_j\nabla u_2\cdot
 \dot{\Delta}_j\nabla\theta{\rm d}x+\int\dot{\Delta}_j u\cdot\dot{\Delta}_j
 \partial_2b{\rm d}x,\\
\begin{aligned}
 \tilde{\mathcal{D}}_H^2
&:= C_2\mu\|\Delta\dot{\Delta}_jb\|_{L^2}^2+C_2\nu\|\nabla\dot{\Delta}_j\theta\|_{L^2}^2
 +\|\dot{\Delta}_j\partial_1 u\|_{L^2}^2+\|\dot{\Delta}_j\partial_2 u\|_{L^2}^2
 +\|\dot{\Delta}_j\partial_2\theta\|_{L^2}^2 \\
&\quad -\|\dot{\Delta}_j\partial_2b\|_{L^2}^2
 -\|\dot{\Delta}_j\nabla \theta\|_{L^2}^2
 -\int\dot{\Delta}_j\nabla\partial_2b_2\cdot\dot{\Delta}_j\nabla\theta{\rm d}x
 +\nu\int\dot{\Delta}_j\nabla\theta\cdot\dot{\Delta}_j\nabla u_2{\rm d}x \\
 &\quad -\int\dot{\Delta}_j\theta\dot{\Delta}_j\partial_2b_2{\rm d}x
  +\nu\int\Delta\dot{\Delta}_jb\cdot\dot{\Delta}_j\partial_2 u{\rm d}x.
 \end{aligned}
\end{gather*}
By Berstein's inequality, for $j\geqslant -1$, we have
\begin{equation}\label{3.17}
 \|\dot{\Delta}_jf\|_{L^2}\lesssim2^{-j}\|\nabla\dot{\Delta}_jf\|_{L^2}\lesssim\|\nabla\dot{\Delta}_jf\|_{L^2}.
\end{equation}
Taking $C_2$ large enough, then using \eqref{3.17} and Young's inequality,
it follows for $j\geqslant -1$ that
\begin{equation}\label{3.18}
\begin{gathered}
 \tilde{\mathcal{E}}_H^2\approx\|(\nabla\dot{\Delta}_ju,\nabla\dot{\Delta}_jb,\nabla\dot{\Delta}_j\theta)\|_{L^2}^2, \\
 \tilde{\mathcal{D}}_H^2\approx\|(\nabla\dot{\Delta}_ju,\Delta\dot{\Delta}_jb,\nabla\dot{\Delta}_j\theta)\|_{L^2}^2, \\
 \tilde{\mathcal{D}}_H^2\gtrsim\|(\nabla\dot{\Delta}_ju,\nabla\dot{\Delta}_jb,\nabla\dot{\Delta}_j\theta)\|_{L^2}\|(\nabla\dot{\Delta}_ju,\Delta\dot{\Delta}_jb,\nabla\dot{\Delta}_j\theta)\|_{L^2}.
\end{gathered}
\end{equation}
Noting that
% \begin{equation}\label{3.19}
\begin{align*}
&\int\nabla\dot{\Delta}_j \mathcal{N}_1\cdot \nabla\dot{\Delta}_ju{\rm d}x
 +\int\nabla\dot{\Delta}_j \mathcal{N}_2\cdot\nabla \dot{\Delta}_jb{\rm d}x
 -\int\nabla\dot{\Delta}_j \mathcal{N}_3\cdot \nabla\dot{\Delta}_j\theta{\rm d}x \\
&=\int\dot{\Delta}_j(\nabla (b\cdot\nabla) b)\cdot\dot{\Delta}_j\nabla u{\rm d}x
 +\int\dot{\Delta}_j(\nabla (b\cdot\nabla) u)\cdot\dot{\Delta}_j\nabla b{\rm d}x
 -\int\dot{\Delta}_j(\nabla (u\cdot\nabla) u)\cdot\dot{\Delta}_j\nabla u{\rm d}x \\
&\quad-\int\dot{\Delta}_j(\nabla (u\cdot\nabla) b)\cdot\dot{\Delta}_j\nabla b{\rm d}x
 -\int\dot{\Delta}_j(\nabla (u\cdot\nabla) \theta)\cdot\dot{\Delta}_j\nabla \theta{\rm d}x
 +\int[\dot{\Delta}_j, b\cdot\nabla]\nabla b\cdot\dot{\Delta}_j\nabla u{\rm d}x \\
 &\quad +\int[\dot{\Delta}_j, b\cdot\nabla]\nabla u\cdot\dot{\Delta}_j\nabla b{\rm d}x
 -\int[\dot{\Delta}_j, u\cdot\nabla]\nabla b\cdot\dot{\Delta}_j\nabla b{\rm d}x
 -\int[\dot{\Delta}_j, u\cdot\nabla]\nabla u\cdot\dot{\Delta}_j\nabla u{\rm d}x \\
&\quad -\int[\dot{\Delta}_j, u\cdot\nabla]\nabla \theta\cdot\dot{\Delta}_j\nabla 
 \theta{\rm d}x
\end{align*}
and
\begin{equation}\label{3.20}
\begin{aligned}
&\int\dot{\Delta}_j\nabla \mathcal{N}_1^2\cdot\dot{\Delta}_j\nabla\theta{\rm d}x
 -\int\dot{\Delta}_j\nabla\mathcal{N}_3\cdot\dot{\Delta}_j\nabla u_2{\rm d}x \\
&=\int\dot{\Delta}_j\nabla (b\cdot\nabla b_2)\cdot\dot{\Delta}_j\nabla \theta{\rm d}x
 -\int\dot{\Delta}_j(\nabla (u\cdot\nabla) u_2)\cdot\dot{\Delta}_j\nabla \theta{\rm d}x\\
&\quad -\int\dot{\Delta}_j(\nabla (u\cdot\nabla) \theta)
 \cdot\dot{\Delta}_j\nabla u_2{\rm d}x
 -\int[\dot{\Delta}_j, u\cdot\nabla]\nabla u_2\cdot\dot{\Delta}_j\nabla \theta{\rm d}x \\
&\quad -\int[\dot{\Delta}_j, u\cdot\nabla]\nabla \theta\cdot\dot{\Delta}_j\nabla u_2{\rm d}x,
\end{aligned}
\end{equation}
then by H\"older's inequality, \eqref{3.16}-\eqref{3.20}, for $j\geqslant -1$, 
one has
\begin{align*}
\frac{\rm d}{{\rm d}t}\tilde{\mathcal{E}}_H^2+\tilde{\mathcal{E}}_H
 \tilde{\mathcal{D}}_H 
&\lesssim\Big{(}\|\dot{\Delta}_j
 (\nabla b\cdot\nabla b,\nabla b\cdot\nabla u,\nabla u
 \cdot\nabla u,\nabla u\cdot\nabla b,
 \nabla u\cdot\nabla \theta)\|_{L^2}\\
&\quad +\|([\dot{\Delta}_j,b\cdot\nabla](\nabla b,\nabla u),[\dot{\Delta}_j,u\cdot\nabla]
 (\nabla b,\nabla u,\nabla\theta))\|_{L^2}
 +\|(\dot{\Delta}_j \mathcal{N}_1,\dot{\Delta}_j \mathcal{N}_2)\|_{L^2}\Big{)}
 \tilde{\mathcal{E}}_H,
\end{align*}
which gives 
 \begin{equation}\label{3.21}
\begin{aligned}
&\tilde{\mathcal{E}}_H(t)+\int_0^t\tilde{\mathcal{D}}_H(s){\rm d}s \\
&\lesssim\tilde{\mathcal{E}}_H(0)+\int_0^t\Big(\|\dot{\Delta}_j
 (\nabla b\cdot\nabla b,\nabla b\cdot\nabla u,\nabla u\cdot\nabla u,
 \nabla u\cdot\nabla b,  \nabla u\cdot\nabla \theta)\|_{L^2} \\
&\quad +\|\Big([\dot{\Delta}_j,b\cdot\nabla](\nabla b,\nabla u),[\dot{\Delta}_j,
 u\cdot\nabla](\nabla b,\nabla u,\nabla\theta)\Big)\|_{L^2}
 +\|(\dot{\Delta}_j \mathcal{N}_1,\dot{\Delta}_j \mathcal{N}_2)\|_{L^2}){\rm d}s.
\end{aligned}
\end{equation}
By using the Calderon-Zygmund inequality, \eqref{2.4}, \eqref{3.18} and summing 
inequality \eqref{3.21} over $j\geqslant -1$, we have
\begin{equation}\label{3.22}
\begin{aligned}
&\|(u,b,\theta)\|_{\tilde{L}^\infty_t(\dot{B}^2_{2,1})}^H
  +\|(u,\theta)\|_{{L}^1_t(\dot{B}^2_{2,1})}^H+\|b\|_{{L}^1_t(\dot{B}^3_{2,1})}^H \\
&\lesssim\|(u_0,b_0,\theta_0)\|_{\dot{B}^2_{2,1}}^H+\int_0^t\Big{(}\|( \mathcal{N}_1, 
 \mathcal{N}_2)\|_{\dot{B}^1_{2,1}} \\
&\quad +\sum_{j\geqslant-1}2^j\|([\dot{\Delta}_j,b\cdot\nabla](\nabla b,\nabla u),
 [\dot{\Delta}_j,u\cdot\nabla](\nabla b,\nabla u,\nabla\theta))\|_{L^2} \\
&\quad +\|(\nabla b\cdot\nabla b,\nabla b\cdot\nabla u,\nabla u\cdot\nabla u,
 \nabla u\cdot\nabla b,\nabla u\cdot\nabla \theta)\|_{\dot{B}^1_{2,1}}\Big{)}{\rm d}s.
\end{aligned}
\end{equation}

Next we turn to estimate the nonlinear terms of \eqref{3.22} in sequence. 
First, by \eqref{2.7} and \eqref{2.1}, we obtain
 \begin{equation}\label{3.23}
\begin{aligned}
&\|(\nabla b\cdot\nabla b,\nabla b\cdot\nabla u,\nabla u
 \cdot\nabla u,\nabla u\cdot\nabla b,\nabla u\cdot\nabla \theta)\|_{\dot{B}^1_{2,1}} \\
&\lesssim\|\nabla b\|_{L^\infty}\|(\nabla b,\nabla u)\|_{\dot{B}^1_{2,1}}
 +\|\nabla u\|_{L^\infty}\|(\nabla b,\nabla u,\nabla\theta)\|_{\dot{B}^1_{2,1}}
 +\|\nabla \theta\|_{L^\infty}\|\nabla u\|_{\dot{B}^1_{2,1}} \\
&\lesssim\|(\nabla b,\nabla u,\nabla\theta)\|_{\dot{B}^1_{2,1}}^2 \\
&\lesssim\|( b, u,\theta)\|_{\dot{B}^2_{2,1}}^2 \\
&\lesssim\mathcal{E}(t)\mathcal{D}(t).
\end{aligned}
\end{equation}
For the term concerning the commutator, we use \eqref{2.11} and \eqref{2.1} to obtain
\begin{equation}\label{3.24}
\begin{aligned}
&\sum_{j\geqslant-1}2^j\|([\dot{\Delta}_j,b\cdot\nabla](\nabla b,\nabla u),
 [\dot{\Delta}_j,u\cdot\nabla](\nabla b,\nabla u,\nabla\theta))\|_{L^2} \\
&\lesssim\|(\nabla b,\nabla u,\nabla\theta)\|_{L^\infty}\|(\nabla b,\nabla u,
 \nabla\theta)\|_{\dot{B}^1_{2,1}} \\
&\lesssim\|(\nabla b,\nabla u,\nabla\theta)\|_{\dot{B}^1_{2,1}}^2 \\
&\lesssim\|( b, u,\theta)\|_{\dot{B}^2_{2,1}}^2 \\
&\lesssim\mathcal{E}(t)\mathcal{D}(t).
\end{aligned}
\end{equation}
For the last term, it follows from \eqref{2.8} and \eqref{2.1} that
\begin{equation}\label{3.25}
 \|( \mathcal{N}_1, \mathcal{N}_2)\|_{\dot{B}^1_{2,1}}
 \lesssim \|( b, u)\|_{\dot{B}^1_{2,1}}\|(\nabla b,\nabla u)\|_{\dot{B}^1_{2,1}}
 \lesssim \|( b, u)\|_{\dot{B}^1_{2,1}}\|( b, u,\theta)\|_{\dot{B}^2_{2,1}} 
 \lesssim \mathcal{E}(t)\mathcal{D}(t).
\end{equation}
Thus, inserting the estimates \eqref{3.23}-\eqref{3.25} into \eqref{3.22}, 
we obtain \eqref{3.14} immediately, and complete the proof.
\end{proof}

\begin{proof}[Proof of Theorem \ref{thm1}]
Here $T^*$ denotes the life-span (maximal existence time) of the local solution 
$(u,b,\theta)$. Next, we need to proof that \eqref{1.3} holds and $T^*=\infty$. 
It infers from \eqref{3.2} and \eqref{3.14} that
\begin{equation}\label{3.26}
 \mathcal{E}(t)+\int_0^t\mathcal{D}(s){\rm d}s\leqslant C_3\mathcal{E}_0+C_4\sup _{0\leqslant s\leqslant t}\mathcal{E}(s)\int_0^t\mathcal{D}(s){\rm d}s,
\end{equation}
where $\mathcal{E}_0$ is defined in \eqref{1.4}.

For any $t\in(0,T^*)$, we have 
\[
 \sup_{0\leqslant s\leqslant t}\mathcal{E}(s)\leqslant\frac{1}{2C_4}.
\]
Then \eqref{3.26} yields
\begin{equation}\label{3.27}
 \mathcal{E}(t)+\frac{1}{2}\int_0^t\mathcal{D}(s){\rm d}s\leqslant C_3\mathcal{E}_0.
\end{equation}
Finally, choosing sufficiently small $\epsilon$, by using the local existence result 
together with a standard continuity argument (sometimes referred to as a bootstrapping 
argument), we have $T^*=\infty$.
This completes the proof. 
\end{proof}

\begin{thebibliography}{99}
\bibitem[21]{MR1107136} 
G.~P. Galdi, M.~Padula; 
\emph{Further results in the nonlinear stability of
 the magnetic {B}\'enard problem}, Mathematical aspects of fluid and plasma
 dynamics ({S}alice {T}erme, 1988), Lecture Notes in Math., vol. 1460,
 Springer, Berlin, 1991, pp.~140--151. MR {1107136}
\bibitem[12]{MR128226} 
S.~Chandrasekhar; 
\emph{Hydrodynamic and hydromagnetic stability}, International Series
 of Monographs on Physics, Clarendon Press, Oxford, 1961. MR {128226}
\bibitem[13]{MR1354312} 
J.-Y. Chemin, N.~Lerner; 
\emph{Flot de champs de vecteurs non lipschitziens
 et \'equations de {N}avier-{S}tokes}, J. Differential Equations,
 \textbf{121} (1995), no.~2, 314--328. MR {1354312}
\bibitem[16]{MR1855277} 
Rapha\"{e}l Danchin; 
\emph{Local theory in critical spaces for compressible
 viscous and heat-conductive gases}, Comm. Partial Differential Equations,
 \textbf{26} (2001), no.~7-8, 1183--1233. MR {1855277}
\bibitem[34]{MR1965452} 
Andrew Majda; 
\emph{Introduction to {PDE}s and waves for the atmosphere and ocean}, 
Courant Lecture Notes in Mathematics, vol.~9, New York University,
Courant Institute of Mathematical Sciences, New York; American Mathematical
Society, Providence, RI, 2003. MR {1965452}
\bibitem[19]{MR2098531} 
P.~G. Drazin, W.~H. Reid; 
\emph{Hydrodynamic stability}, second ed.,
 Cambridge Mathematical Library, Cambridge University Press, Cambridge, 2004,
 With a foreword by John Miles. MR {2098531}
\bibitem[1]{MR2418123} 
Hammadi Abidi, Marius Paicu; 
\emph{Global existence for the
 magnetohydrodynamic system in critical spaces}, Proc. Roy. Soc. Edinburgh
 Sect. A ,\textbf{138} (2008), no.~3, 447--476. MR {2418123}
\bibitem[3]{MR2768550} 
Hajer Bahouri, Jean-Yves Chemin, Rapha\"{e}l Danchin; 
\emph{Fourier  analysis and nonlinear partial differential equations}, 
Grundlehren der  mathematischen Wissenschaften [Fundamental Principles of Mathematical
 Sciences], vol. 343, Springer, Heidelberg, 2011. MR {2768550}
\bibitem[33]{MR3168121} 
Fanghua Lin, Ping Zhang; 
\emph{Global small solutions to an {MHD}-type
 system: the three-dimensional case}, Comm. Pure Appl. Math., \textbf{67}
 (2014), no.~4, 531--580. MR {3168121}
\bibitem[35]{MR3210038} 
Xiaoxia Ren, Jiahong Wu, Zhaoyin Xiang, Zhifei Zhang; 
\emph{Global  existence and decay of smooth solution for the 2-{D} {MHD} equations
 without  magnetic diffusion}, J. Funct. Anal. \textbf{267} (2014), no.~2, 503--541.
 MR {3210038}
\bibitem[29]{MR3245016} 
Alexander Kiselev, Vladimir \v{S}ver\'{a}k; 
\emph{Small scale creation for  solutions of the incompressible two-dimensional 
{E}uler equation}, Ann. of  Math., (2) \textbf{180} (2014), no.~3, 1205--1220. MR {3245016}
\bibitem[18]{MR3293735} 
Sergey~A. Denisov; 
\emph{Double exponential growth of the vorticity gradient
 for the two-dimensional {E}uler equation}, Proc. Amer. Math. Soc.,
 \textbf{143} (2015), no.~3, 1199--1210. MR {3293735}
\bibitem[8]{MR3363411} 
J.~Boussinesq; 
\emph{Th\'eorie des ondes et des remous qui se propagent le
 long d'un canal rectangulaire horizontal, en communiquant au liquide contenu
 dans ce canal des vitesses sensiblement pareilles de la surface au fond}, J.
 Math. Pures Appl., (2) \textbf{17} (1872), 55--108. MR {3363411}
\bibitem[32]{MR3377532} 
Fanghua Lin, Li~Xu, Ping Zhang; 
\emph{Global small solutions of 2-{D}
 incompressible {MHD} system}, J. Differential Equations, \textbf{259} (2015),
 no.~10, 5440--5485. MR {3377532}
\bibitem[15]{MR3412279} 
Jianfeng Cheng, Lili Du; 
\emph{On two-dimensional magnetic {B}\'enard
 problem with mixed partial viscosity}, J. Math. Fluid Mech., \textbf{17}
 (2015), no.~4, 769--797. MR {3412279}
\bibitem[46]{MR3448784} 
Ting Zhang; 
\emph{Global solutions to the 2{D} viscous, non-resistive {MHD}
system with large background magnetic field}, J. Differential Equations,
\textbf{260} (2016), no.~6, 5450--5480. MR {3448784}
\bibitem[20]{MR346289} 
G.~Duvaut, J.-L. Lions; 
\emph{In\'equations en thermo\'elasticit\'e et
 magn\'etohydrodynamique}, Arch. Rational Mech. Anal., \textbf{46} (1972),
 241--279. MR {346289}
\bibitem[42]{MR3475939} 
Xiaoqian Xu; 
\emph{Fast growth of the vorticity gradient in symmetric smooth
 domains for 2{D} incompressible ideal flow}, J. Math. Anal. Appl.,
 \textbf{439} (2016), no.~2, 594--607. MR {3475939}
\bibitem[2]{MR3666563} 
Hammadi Abidi, Ping Zhang; 
\emph{On the global solution of a 3-{D} {MHD}
 system with initial data near equilibrium}, Comm. Pure Appl. Math.,
 \textbf{70} (2017), no.~8, 1509--1561. MR {3666563}
\bibitem[40]{MR3678491} 
Dongyi Wei, Zhifei Zhang; 
\emph{Global well-posedness of the {MHD} equations
 in a homogeneous magnetic field}, Anal. PDE, \textbf{10} (2017), no.~6,
 1361--1406. MR {3678491}
\bibitem[5]{MR3710679} 
Dongfen Bian and Jitao Liu; 
\emph{Initial-boundary value problem to 2{D}
 {B}oussinesq equations for {MHD} convection with stratification effects}, J.
 Differential Equations, \textbf{263} (2017), no.~12, 8074--8101. MR {3710679}
\bibitem[22]{MR3738384} 
Ling-Bing He, Li~Xu, Pin Yu; 
\emph{On global dynamics of three dimensional
 magnetohydrodynamics: nonlinear stability of {A}lfv\'en waves}, Ann. PDE,
 \textbf{4} (2018), no.~1, Paper No. 5, 105. MR {3738384}
\bibitem[9]{MR3780143} 
Yuan Cai, Zhen Lei; 
\emph{Global well-posedness of the incompressible
 magnetohydrodynamics}, Arch. Ration. Mech. Anal., \textbf{228} (2018), no.~3,
 969--993. MR {3780143}
\bibitem[47]{MR3843627} 
Yi~Zhou, Yi~Zhu; 
\emph{Global classical solutions of 2{D} {MHD} system with
 only magnetic diffusion on periodic domain}, J. Math. Phys., \textbf{59}
 (2018), no.~8, 081505, 12. MR {3843627}
\bibitem[17]{MR3851055} 
Wen Deng, Ping Zhang; 
\emph{Large time behavior of solutions to 3-{D} {MHD}
 system with initial data near equilibrium}, Arch. Ration. Mech. Anal.,
 \textbf{230} (2018), no.~3, 1017--1102. MR {3851055}
\bibitem[45]{MR3887133} 
Jianwen Zhang, Mingyu Zhang; 
\emph{On the infinite {P}randtl number limit in
 two-dimensional magneto-convection}, Nonlinear Anal. Real World Appl.,
 \textbf{46} (2019), 313--334. MR {3887133}
\bibitem[37]{MR3977115} 
Haifeng Shang; 
\emph{Global regularity results for the 2{D} magnetic
 {B}\'enard problem with fractional dissipation}, J. Math. Fluid Mech.,
 \textbf{21} (2019), no.~3, Paper No. 39, 35. MR {3977115}
\bibitem[24]{MR3980230} 
Ruihong Ji, Dan Li, Youhua Wei,  Jiahong Wu; 
\emph{Stability of hydrostatic
 equilibrium to the 2{D} {B}oussinesq systems with partial dissipation}, Appl.
 Math. Lett., \textbf{98} (2019), 392--397. MR {3980230}
\bibitem[7]{MR4046053} 
Dongfen Bian and Xueke Pu; 
\emph{Global smooth axisymmetic solutions of the
 {B}oussinesq equations for magnetohydrodynamics convection}, 
J. Math. Fluid  Mech., \textbf{22} (2020), no.~1, Paper No. 12, 13. MR {4046053}
\bibitem[25]{MR4134456} 
Fei Jiang, Song Jiang; 
\emph{On inhibition of thermal convection instability
 by a magnetic field under zero resistivity}, J. Math. Pures Appl., (9)
 \textbf{141} (2020), 220--265. MR {4134456}
\bibitem[41]{MR4199961} 
Dongyi Wei, Zhifei Zhang; 
\emph{Global well-posedness for the 2-{D} {MHD} equations with
 magnetic diffusion}, Commun. Math. Res., \textbf{36} (2020), no.~4, 377--389.
 MR {4199961}
\bibitem[26]{MR4339533} 
Fei Jiang, Song Jiang; 
\emph{Asymptotic behaviors of global solutions to the two-dimensional
 non-resistive {MHD} equations with large initial perturbations}, Adv. Math.,
 \textbf{393} (2021), Paper No. 108084, 79. MR {4339533}
\bibitem[14]{MR4372923} 
Wenji Chen, Zhifei Zhang, Jianfeng Zhou; 
\emph{Global well-posedness for
 the 3-{D} {MHD} equations with partial diffusion in the periodic domain},
 Sci. China Math., \textbf{65} (2022), no.~2, 309--318. MR {4372923}
\bibitem[39]{MR44301} 
W.~B. Thompson; 
\emph{Thermal convection in a magnetic field}, Proc. {S}econd
 {C}anadian {M}ath. {C}ongress, {V}ancouver, 1949, Univ. Toronto Press,
 Toronto, ON, 1951, pp.~196--205. MR {44301}
\bibitem[6]{MR4484087} 
Dongfen Bian, Jingjing Mao, and Xueke Pu; 
\emph{Stability of hydrostatic
 equilibrium for the 2{D} {BMHD} system with partial dissipation}, Commun.
 Pure Appl. Anal., \textbf{21} (2022), no.~10, 3441--3462. MR {4484087}
\bibitem[43]{MR4496608} 
Wei~Kui Ye, Zhao~Yang Yin; 
\emph{Global well-posedness for the non-viscous
 {MHD} equations with magnetic diffusion in critical {B}esov spaces}, Acta
 Math. Sin. (Engl. Ser.), \textbf{38} (2022), no.~9, 1493--1511. MR {4496608}
\bibitem[38]{MR44996} 
W.~B. Thompson; 
\emph{Thermal convection in a magnetic field}, Philos. Mag., (7)
 \textbf{42} (1951), 1417--1432. MR {44996}
\bibitem[23]{MR4500250} 
Wenting Huang; 
\emph{Stability and exponential decay of the magnetic
 {B}\'enard system with horizontal dissipation}, J. Math. Anal. Appl.,
 \textbf{519} (2023), no.~1, Paper No. 126767, 25. MR {4500250}
\bibitem[31]{MR4500891} 
Xinliang Li, Zhong Tan, and Saiguo Xu; 
\emph{Global existence and decay  estimates of solutions to the 
{MHD}-{B}oussinesq system with stratification
 effects}, Nonlinearity, \textbf{35} (2022), no.~12, 6067--6097. MR {4500891}
\bibitem[28]{MR4505194} 
Fei Jiang and Yanjin Wang; 
\emph{Nonlinear stability of the inviscid magnetic
 {B}\'enard problem}, J. Math. Fluid Mech., \textbf{24} (2022), no.~4, Paper
 No. 110, 18. MR {4505194}
\bibitem[44]{MR4624442} 
Xiaoping Zhai; 
\emph{Stability for the 2{D} incompressible {MHD} equations with
 only magnetic diffusion}, J. Differential Equations, \textbf{374} (2023),
 267--278. MR {4624442}
\bibitem[27]{MR4641656} 
Fei Jiang, Song Jiang; 
\emph{On magnetic inhibition theory in 3{D} non-resistive
 magnetohydrodynamic fluids: global existence of large solutions}, Arch.
 Ration. Mech. Anal., \textbf{247} (2023), no.~5, Paper No. 96, 35.
 MR {4641656}
\bibitem[30]{MR4790886} 
Suhua Lai, Linxuan Shen, Xia Ye, Xiaokui Zhao; 
\emph{The stabilizing effect  of temperature and magnetic field on a 2{D} 
magnetic {B}\'enard fluids}, J.
 Differential Equations, \textbf{409} (2024), 851--880. MR {4790886}
\bibitem[10]{MR53708} 
S.~Chandrasekhar; 
\emph{On the inhibition of convection by a magnetic field},
 Philos. Mag., (7) \textbf{43} (1952), 501--532. MR {53708}
\bibitem[11]{MR64593} 
S.~Chandrasekhar; 
\emph{On the inhibition of convection by a magnetic field. {II}},
 Philos. Mag., (7) \textbf{45} (1954), 1177--1191. MR {64593}
\bibitem[36]{MR716200} 
Michel Sermange, Roger Temam; 
\emph{Some mathematical questions related to
 the {MHD} equations}, Comm. Pure Appl. Math. \textbf{36} (1983), no.~5,
 635--664. MR {716200}
\bibitem[4]{MR920153} 
C.~Bardos, C.~Sulem, P.-L. Sulem; 
\emph{Longtime dynamics of a conductive
 fluid in the presence of a strong magnetic field}, Trans. Amer. Math. Soc.
 \textbf{305} (1988), no.~1, 175--191. MR {920153}
\end{thebibliography}
\end{document}
